\chapter{Introduction}
\section{Background}
Business and public-sector workflows depend on information received in semi-structured or unstructured documents. A clerk opens a PDF, finds the supplier name, invoice number, date, tax amount and total, and types them into a business application. Repeated across hundreds of documents, this consumes staff time and introduces errors.

OCR converts visual text into machine-readable text, but OCR alone does not solve the problem. A practical system must also locate the relevant fields, normalise values, apply business rules, preserve the source document, track corrections and deliver approved data to the target system.

\section{Problem Statement}
There is a gap between incoming digital documents and structured business records. A basic OCR tool returns text without knowing which text is the invoice number, or whether the total agrees with the subtotal and tax. A useful pipeline must connect acquisition, recognition, extraction, validation, human review and controlled ERP delivery in a traceable workflow.

\section{Aim}
To design and implement a secure, modular document-processing system that extracts structured information from French and English invoices and prepares validated records for ERP ingestion.

\section{Specific Objectives}
\begin{enumerate}
\item Provide secure document upload and lifecycle tracking.
\item Preprocess scanned images to improve text readability.
\item Develop and evaluate a trainable OCR model for French and English text.
\item Extract key invoice fields using configurable rules and layout information.
\item Compute field-level confidence and identify records that need human review.
\item Validate fields using format checks, required-field checks, duplicate detection and arithmetic rules.
\item Provide a dashboard to review, correct, approve, reject and track documents.
\item Store metadata, extracted values, validation outcomes and audit events in MySQL.
\item Expose controlled interfaces to an ERP adapter or a sandbox ERP endpoint.
\item Package the system with Docker and document a Kubernetes path.
\end{enumerate}

\section{Research and Engineering Questions}
\begin{itemize}
\item How much does image preprocessing improve recognition on the selected dataset?
\item How does the custom OCR model compare with an established baseline engine?
\item Which combination of confidence, validation rules and human review balances automation and correctness?
\item How can documents be processed reliably without losing traceability?
\item How can ERP-specific code be isolated from document-processing logic?
\end{itemize}

\section{Expected Contribution}
The contribution is an integrated, measurable prototype, not an unsupported claim that a custom OCR model beats existing tools. The work combines a trainable recogniser, domain-specific extraction and validation, human review, auditability and a replaceable ERP adapter. The comparison experiment will show where the custom component works well and where it needs improvement.

\begin{infobox}[Project name]
The working name \projname{} comes from the naming shortlist in the planning phase. It is a suggestion, not a verified trademark or domain. The name is defined in one place (\texttt{\textbackslash projname}) so it can be changed everywhere at once.
\end{infobox}

\chapter{Scope, Users and Assumptions}
\section{Initial Scope}
The first version handles \textbf{invoices} only (supplier invoices received as PDF, PNG or JPEG), in French and English. Purchase orders, receipts, forms and administrative requests are extensions for later versions. The ERP stage uses a sandbox API or a clearly labelled mock adapter until the target ERP and its interface are confirmed.

\section{Out of Scope for the First Prototype}
\begin{itemize}
\item Fully autonomous posting of every document without human review.
\item Support for every document type, language, handwriting style or layout.
\item A production connection to an ERP whose API, schema and authentication are not yet known.
\item Claims of legal or accounting compliance without review against applicable rules.
\item Kubernetes operations before the containerised application is stable and tested.
\end{itemize}

\section{Stakeholders and Roles}
{\tblstart\small
\begin{longtable}{L{3.0cm}L{4.4cm}L{6.6cm}}
\hdrrow \hdr{Role} & \hdr{Responsibility} & \hdr{Permissions / use}\\ \endhead
Administrator & Manage users, configuration, thresholds and integration settings. & Create or deactivate users, assign roles, view system status, configure rules.\\
Document operator & Upload and monitor incoming documents. & Upload, view status of permitted records.\\
Reviewer / accountant & Check extracted fields and resolve uncertainty. & Correct fields, add notes, approve or reject records.\\
Integration service account & Transfer approved records to the ERP. & Restricted API credential; cannot administer users or change source documents.\\
Supervisor / auditor & Assess operation and traceability. & Read authorised logs, metrics and audit history.\\
\hline
\end{longtable}}

\section{Assumptions and Dependencies}
\begin{itemize}
\item A labelled dataset of representative French and English invoices can be obtained lawfully or created (including synthetic invoices).
\item Training and evaluation data can be separated to avoid leakage.
\item The development machine runs a stable Python release; a GPU (or a cloud notebook) is optional for the prototype but useful for training.
\item MySQL is available locally during early development.
\item The target ERP and its integration method will be chosen before any production posting.
\item Test documents contain no real sensitive data unless authorised and protected.
\end{itemize}

\chapter{Requirements Specification}
Requirements use stable identifiers so that each can be traced to a design element and a test. The full Phase 1 specification, with acceptance criteria and a traceability matrix, is in the companion document.

\section{Functional Requirements}
{\tblstart\small
\begin{longtable}{L{1.4cm}L{11.1cm}L{1.4cm}}
\hdrrow \hdr{ID} & \hdr{Requirement} & \hdr{Priority}\\ \endhead
FR-01 & A permitted user can sign in and sign out. & \must\\
FR-02 & The system enforces role-based access to administrative, operational and review functions. & \must\\
FR-03 & A user can upload a PDF, PNG or JPEG document within configured size and type limits. & \must\\
FR-04 & The system checks file type, size and basic integrity before accepting a document. & \must\\
FR-05 & The original file is stored in protected storage, linked to a unique document identifier and a SHA-256 hash. & \must\\
FR-06 & The system creates a processing job and records its status and timestamps. & \must\\
FR-07 & The preprocessing module normalises pages and applies suitable image enhancement. & \must\\
FR-08 & The OCR component returns recognised text and confidence information where supported. & \must\\
FR-09 & The extraction component maps recognised content to the configured invoice fields. & \must\\
FR-10 & The validation component checks required fields, formats, arithmetic consistency and duplicate indicators. & \must\\
FR-11 & Low-confidence or invalid records are marked for human review. & \must\\
FR-12 & An authorised reviewer can correct fields, add a reason, and approve or reject a record. & \must\\
FR-13 & Every material correction, approval, rejection and integration attempt is recorded in an audit trail. & \must\\
FR-14 & Only approved records can be submitted to the ERP adapter. & \must\\
FR-15 & The system records ERP submission status, external reference and errors without duplicating a successful submission. & \must\\
FR-16 & Users can search and filter documents by status, date, type or identifier. & \should\\
FR-17 & Administrators can configure validation rules and review thresholds. & \should\\
FR-18 & The system exposes health endpoints and basic operational metrics. & \should\\
FR-19 & The dashboard shows processing counts and the failure and review queues. & \should\\
FR-20 & Approved structured data can be exported (JSON/CSV) when ERP connectivity is unavailable. & \could\\
FR-21 & A failed or rejected document can be resubmitted for reprocessing as a new processing run. & \should\\
FR-22 & The review screen shows the source page beside the extracted fields and highlights the evidence for each field. & \should\\
\hline
\end{longtable}}

\section{Non-Functional Requirements}
{\tblstart\small
\begin{longtable}{L{1.7cm}L{2.5cm}L{9.7cm}}
\hdrrow \hdr{ID} & \hdr{Category} & \hdr{Requirement and verification approach}\\ \endhead
NFR-01 & Security & Authenticate users, authorise every protected operation, keep secrets out of source code, restrict stored files. Verified by access-control tests.\\
NFR-02 & Privacy & Collect only necessary data, define retention, use synthetic or authorised test documents.\\
NFR-03 & Reliability & A failed processing step yields a visible failure state and a recoverable record, never a silently lost job.\\
NFR-04 & Traceability & Every document and processing run has stable identifiers and timestamps.\\
NFR-05 & Maintainability & Separate API, domain logic, OCR, validation, storage and ERP adapter responsibilities.\\
NFR-06 & Portability & Reproducible setup through documented steps and Docker images.\\
NFR-07 & Performance & Measure upload, preprocessing, OCR, extraction and end-to-end times on a defined test set; fix targets after a baseline run.\\
NFR-08 & Usability & Reviewers can find low-confidence fields and understand validation errors without reading logs; tested with representative users.\\
NFR-09 & Scalability & Long-running OCR is separable from interactive API requests and can later scale independently.\\
NFR-10 & Observability & Structured logs, health status, failure categories and durations, without logging passwords or full document content.\\
\hline
\end{longtable}}

\section{Business Rules}
\begin{enumerate}
\item A document cannot be submitted to an ERP until it has passed required validations and been approved by an authorised reviewer.
\item A low-confidence field must be visible to the reviewer; confidence is an indicator, not proof of correctness.
\item Duplicate detection flags likely duplicates for review; it never deletes a document automatically.
\item ERP requests use an idempotency key or an equivalent duplicate-prevention mechanism.
\item Corrections preserve the original extracted value and record the corrected value, actor, time and reason.
\item Failed integration attempts can be retried without creating duplicate business records.
\end{enumerate}

\section{Provisional Success Metrics}
The targets below are \textbf{planning values}, chosen to give the experiments something concrete to aim at. They must be revised after the baseline measurement and must never be reported as results.

{\tblstart\small
\begin{longtable}{L{4.3cm}L{4.4cm}L{4.4cm}}
\hdrrow \hdr{Metric} & \hdr{Provisional target} & \hdr{How measured}\\ \endhead
Character error rate (CER) & Lower than the baseline engine on the same test set; absolute target set after baseline & Edit distance vs.\ verified transcription\\
Word error rate (WER) & As above & Word-level edit distance\\
Field F1 (invoice no., date, total) & $\geq 0.85$ after review-free extraction, to be revised & Exact/normalised match on held-out invoices\\
Records needing review & Reported, not targeted; aim to reduce over time & Share of records flagged\\
Processing time per page & Reported as median and 95th percentile & Timed end-to-end runs\\
Duplicate ERP submissions & 0 in all retry tests & Sandbox integration tests\\
\hline
\end{longtable}}
